\addtocounter{page}{4}

\chapter*{РЕФЕРАТ}
\addcontentsline{toc}{chapter}{РЕФЕРАТ}

Ключевые слова:
\begin{itemize}
    \item рекомендательная система;
    \item коллаборативная фильтрация;
    \item контентная фильтрация;
    \item регрессия;
    \item Android.
\end{itemize}

Цель работы: разработка системы рекомендаций рецептов блюд на основе предпочтений и возможностей пользователя.
В задачи работы входит выбор алгоритма рекомендаций, его проектирование и реализация.
Рассмотрено $8$ методов формирования рекомендаций.
Для каждого из них приведено подробное описание.
Далее рассмотрены особенности каждого из методов и проведено их сравнение по сформулированным критериям.
Также представлена реализация преобразования текста рецептов блюд в вектора вещественных чисел.
Перед определением требований к разрабатываемой программе проанализированы и сравнены существующие аналоги.
Выбранный алгоритм рекомендаций -- линейная регрессия с регуляризацией после уменьшения размерности векторов рецептов.
В ходе проектирования были представлены схемы алгоритма работы рекомендательной системы, диаграмма прецедентов, диаграмма сущность-связь и диаграмма коммуникаций для выбранного архитектурного шаблона проектирования <<MVI>>.
Исследована зависимость метрик качества рекомендательной системы от выбора алгоритма рекомендаций.
В результате даны достаточные сведения для реализации рекомендательной системы рецептов блюд.
В данной работе не решается проблема холодного старта.
Это значит, что все рассмотренные методов опираются лишь на данные о предпочтениях пользователя, полученные в ходе его взаимодействия с системой, а не на знания о данном пользователе.